\section{A commuting operator model}
\label{sec:r3-normal-form}

All generalized polarizations act as vector fields on a single
family of pointed jet determinants.  An invertible source jet can
be reparametrized to $t$; the other jet coefficients then become
commuting polarizations from that source row.  Polynomial dependence
extends this description across the singular jets.  The resulting
normal form makes specialization in the number of alphabets exact.

\begin{lemma}[Normal ordering]\label{lem:r2-normal-order}
Let $V$ be a finite-dimensional row space over a field of
characteristic zero.  Suppose $\mathcal M(V)$ is generated from alternating harmonic
polynomials by partial derivatives and generalized polarizations.
Let $\mathcal B(V)$ be the span of the words using only generalized
polarizations on those generators.  Then
\[
 \mathcal M(V)=K[\partial]\mathcal B(V),
 \qquad
 \mathcal M(V)^{\sgn}=\mathcal B(V).
\]
In particular, the alternating part of $\Pol$ is generated from the
ordinary Vandermonde by generalized polarizations alone.
\end{lemma}

\begin{proof}
For $E=\sum_i x_i^{(a)}(\partial_{x_i^{(b)}})^r$,
\[
 [E,\partial_{x_j^{(c)}}]
       =-\delta_{ac}(\partial_{x_j^{(b)}})^r.
\]
More generally, the commutator of $E$ with any constant coefficient
differential operator is again such an operator.  Commuting partial
derivatives to the left, by induction on the number of polarization
operators and on their positions in a word, proves the first identity.
Every element of $\mathcal B(V)$ is alternating, because the
polarizations commute with $S_n$.

Apply the sign projector to an expression $\sum_\ell D_\ell g_\ell$,
where $g_\ell\in\mathcal B(V)$.  Since each $g_\ell$ is alternating,
the result is $\sum_\ell\Av_{S_n}(D_\ell)g_\ell$.  The positive-degree
invariant differential operators annihilate every diagonal harmonic.
Only the constant term of each averaged differential operator remains.
This proves the second identity.
\end{proof}


Put $N=\binom n2$, write $x_i=x_i^{(1)}$, and normalize
\[
 \Delta_n=\det(x_i^j)_{\substack{1\leq i\leq n\\0\leq j<n}}.
\]
All fields in this section have characteristic zero.  For $a\geq2$ put
$Y_{a,s}=\sum_i x_i^{(a)}\partial_{x_i}^{s}$.
The operators $Y_{a,s}$ and the individual $\partial_{x_i}$ commute.
Let $\mathcal B_{n,k}$ be the span of all words in generalized
polarizations applied to $\Delta_n$.

\begin{theorem}[Commuting normal form]\label{thm:r3-normal-form}
For every $n,k\geq1$,
\begin{align}
 \mathcal B_{n,k}
   &=K[Y_{a,s}:2\leq a\leq k,\ 1\leq s<n]\,\Delta_n,
       \label{eq:r3-pure-normal-form}\\
 \Pol_{n,k}
   &=K[\partial_{x_1},\ldots,\partial_{x_n},
           Y_{a,s}:2\leq a\leq k,\ 1\leq s<n]\,\Delta_n.
       \label{eq:r3-full-normal-form}
\end{align}
Moreover, $\Pol_{n,k}^{\sgn}=\mathcal B_{n,k}$.
Thus one source row and a commuting family of operators suffice
for the module defined using arbitrary generalized polarizations.
\end{theorem}

\begin{proof}
The assertion for $n=1$ is immediate, so assume $n\geq2$.
Let $\mathfrak t=tK[t]/(t^n)$, and for
$f=(f_1,\ldots,f_k)\in\mathfrak t^k$ define
\begin{equation}\label{eq:r3-jet-determinant}
 G_f(X)=\left(\prod_{j=0}^{n-1}j!\right)
 \det\left([t^j]\exp\left(\sum_{a=1}^k
                              x_i^{(a)}f_a(t)\right)\right)_{i,j}.
\end{equation}
Only coefficients of order less than $n$ are used.  Thus this is a
polynomial in $X$ and in the finitely many coefficients of $f$,
of $X$-degree at most $N$.
Let $J$ be the span of its values as $f$ varies.

The one-particle identity
\[
 x^{(a)}\partial_{x^{(b)}}^s e^{\sum_c x^{(c)}f_c(t)}
      =x^{(a)}f_b(t)^s e^{\sum_c x^{(c)}f_c(t)}
\]
and differentiation of a determinant give
\begin{equation}\label{eq:r3-jet-vector-field}
 E_{a,b}^{(s)}G_f
     =\left.\frac{d}{d\epsilon}
               G_{f+\epsilon e_a f_b^s}\right|_{\epsilon=0}.
\end{equation}
The product $f_b^s$ is taken modulo $t^n$ and has zero constant
term.  A derivative of a polynomial family lies in the span of its
values, by polynomial interpolation over an infinite field.
Consequently $J$ is stable under every generalized polarization,
including those with output in the source row.
It contains $G_{(t,0,\ldots,0)}=\Delta_n$, and hence
$\mathcal B_{n,k}\subseteq J$.

Suppose that $f_1'(0)\ne0$, and let $g\in\mathfrak t$ be its
compositional inverse modulo $t^n$.
Composition by $g$ acts on coefficients of degrees $0,\ldots,n-1$
through a triangular matrix with diagonal
$1,g'(0),\ldots,g'(0)^{n-1}$.  Therefore
\begin{equation}\label{eq:r3-reparametrization}
 G_{f\circ g}=g'(0)^N G_f.
\end{equation}
Since $f_1\circ g=t$, write
$f_a\circ g=\sum_{s=1}^{n-1}c_{a,s}t^s$ for $a\geq2$.
Expanding the commuting operator exponential gives
\begin{equation}\label{eq:r3-normalized-jet}
 G_{f\circ g}
   =\exp\left(\sum_{a=2}^k\sum_{s=1}^{n-1}
                       c_{a,s}Y_{a,s}\right)\Delta_n.
\end{equation}
The exponential is finite on $\Delta_n$, since every $Y_{a,s}$
lowers its source-row degree by $s$.
Thus $G_f$ belongs to the right side of
\eqref{eq:r3-pure-normal-form} whenever $f_1'(0)\ne0$.
That condition defines a dense open subset of the jet parameter
space.  A polynomial family taking values in a fixed linear
subspace on a dense subset takes values in that subspace everywhere:
apply any linear functional vanishing on the subspace.
It follows that $J$ is contained in the right side of
\eqref{eq:r3-pure-normal-form}.  The reverse inclusion is immediate,
which proves the first identity.

Lemma~\ref{lem:r2-normal-order} gives
$\Pol_{n,k}=K[\partial_{x_i^{(a)}}]\mathcal B_{n,k}$.
Substitute \eqref{eq:r3-pure-normal-form} and move every
derivative in an auxiliary row to the right, using
\[
 [\partial_{x_i^{(a)}},Y_{b,s}]
           =\delta_{ab}\partial_{x_i}^{s}\qquad(a,b\geq2).
\]
Those derivatives kill $\Delta_n$ when they reach it.
The remaining operators are exactly those in
\eqref{eq:r3-full-normal-form}.

The same normal-ordering lemma identifies the alternating part
with $\mathcal B_{n,k}$.
\end{proof}

Positive-order polarizations preserve the pointed condition
$f_a(0)=0$ in \eqref{eq:r3-jet-vector-field}.  Individual
translations of the source variables encode partial derivatives
and give the following finite generating family.

\begin{corollary}[A finite generating determinant]
\label{cor:r3-generating-determinant}
The coefficient span, in the parameters $u_i$ and $c_{a,s}$, of
\begin{equation}\label{eq:r3-generating-determinant}
 \left(\prod_{j=0}^{n-1}j!\right)
 \det\left([t^j]\exp\left((x_i+u_i)t+
         \sum_{a=2}^k x_i^{(a)}
                  \sum_{s=1}^{n-1}c_{a,s}t^s\right)\right)_{i,j}
\end{equation}
is exactly $\Pol_{n,k}$.
Without the parameters $u_i$, the coefficient span is its alternating
part.  A spanning set consists of the words
\[
 \partial_x^\alpha\prod_{a,s}Y_{a,s}^{m_{a,s}}\Delta_n,
 \qquad |\alpha|+\sum_{a,s}s\,m_{a,s}\leq N.
\]
\end{corollary}

\begin{proof}
The determinant is
$\exp(\sum_i u_i\partial_{x_i}+\sum_{a,s}c_{a,s}Y_{a,s})\Delta_n$.
Use Theorem~\ref{thm:r3-normal-form} and compare coefficients.
The displayed inequality is the available source-row degree.
\end{proof}

\begin{corollary}[Exact specialization in the number of rows]
\label{cor:r3-exact-specialization}
If $r\geq k$, setting rows $k+1,\ldots,r$ equal to zero maps
$\Pol_{n,r}$ onto $\Pol_{n,k}$.  Moreover,
\[
 \Pol_{n,r}\cap K[x_i^{(a)}:a\leq k]=\Pol_{n,k}.
\]
The stable row character specializes to the actual finite-row module
for every $k$.
\end{corollary}

\begin{proof}
In the commuting normal form, the auxiliary-row degree of a word is
exactly the number of its corresponding $Y$ factors.  Setting a row
to zero kills precisely the words using a polarization into that row.
All remaining words belong to $\Pol_{n,k}$, and every such word
already occurs in $\Pol_{n,r}$.
If a polynomial is independent of the discarded rows, restriction
fixes it, proving the intersection assertion.
For a polynomial row representation, the zero-weight subspace in
the discarded rows is the usual specialization of its Schur character.
\end{proof}

Haiman's bivariate operator theorem
\cite[Theorem~4.2]{Haiman} describes the full bivariate harmonic
space using the commuting operators in
\eqref{eq:r3-full-normal-form}.  Bergeron's multivariate module
\cite[Section~2.1]{Bergeron} is defined by closure under derivatives
in all variables and polarizations between arbitrary pairs of rows.
Theorem~\ref{thm:r3-normal-form} identifies these operator
descriptions in every number of rows and makes their finite-row
specialization exact.

